% !TEX root = ../main.tex

% 分析与讨论	
\begin{table}
	\renewcommand\arraystretch{1.7}
	\begin{tabularx}{\textwidth}{|X|X|X|X|}
		\hline
		Major：& Physics &Grade：& 2022\\
		\hline
		Name： & 杨舒云 \& 戴鹏辉 & Student number：& 22344020 \& 22344016\\
		\hline
		Date：& 2024/9/25 & Score： &\\
		\hline
	\end{tabularx}
\end{table}
\section{ET2-1 Welding and Debugging of Bluetooth Speakers \\ Analysis \& Discussion}


%---------------------------------------------------------------------
% 数据处理
\subsection{Data Processing}

% 分析
\subsubsection{Analysis}
In this part, we will solve the requested problems .
\begin{enumerate}
	\item Analyze the function of the circled circuit in \cref{fig:question}.
	
	\begin{figure}[htbp]
		\centering
		\includegraphics[width=0.7\textwidth]{ET2_1_question.png}
		\caption{The circled circuit}
		\label{fig:question}
	\end{figure}
	
	The circuit is powered through a USB cable, where Q1, R0, S1, and C0 form an electronic switch to replace a common mechanical switch. S1 is a mercury switch, which relies on the flow of mercury to establish or break the connection. When the speaker is placed upright, the mercury switch conducts. If the speaker is tilted, the mercury switch breaks. This principle is used to control the on-off state of the circuit. Q1 is an N-channel MOSFET, which requires a positive voltage at its gate to conduct. When the mercury switch is on, the gate receives positive voltage, and Q1 conducts, providing power to the amplifier circuit. When the mercury switch is off, the gate is grounded through resistor R0, causing Q1 to turn off, thereby cutting off power to the amplifier circuit. Capacitor Co is used for switch debounce. Additionally, the operating current of the mercury switch is small; a 3 mm mercury switch cannot handle more than 300 mA of current. The speaker, however, requires approximately 1 A of operating current, hence a MOSFET is needed for control.
	
	More information:
	\begin{itemize}
		\item MOSFET（Metal-Oxide-Semiconductor Field-Effect Transistor）是一种场效应晶体管，通常用于开关和放大应用。它具有三个主要电极：栅极（Gate, G）、漏极（Drain, D）、源极（Source, S）。本电路中使用的是P沟道MOSFET（型号为2302），它的工作原理是当栅极电压比源极电压更低时，MOSFET导通，从源极到漏极的电流被允许通过；而当栅极电压和源极电压相等或栅极电压比源极电压高时，MOSFET关闭，从而阻止电流流动。
		
		\item 在电路中，电源通过Micro USB接口输入，并提供VDD电压。电容C0（10μF）作为旁路电容，用于滤波，降低电源电压中的噪声，提供稳定的直流电压给后续电路。PMOS管Q1（型号为2302）是整个电路的核心开关元件，它与机械开关S1一起用于控制电源是否为后续电路供电。当S1闭合时，栅极通过电阻R0（20kΩ）接地，形成一个负栅源电压，使Q1导通，从而使VDD电压能够传递给其他部分的电路；而当S1断开时，栅极通过R0与源极电压趋于同电位（无负压），使Q1截止，切断电源。电阻R0的作用是在开关S1打开时，将Q1的栅极电压拉高至接近源极电压（即VDD），使得Q1能够正确关闭。当S1闭合时，R0限制了从栅极到地的电流，确保栅极电压足够低于源极，使Q1导通。另外，PMOS管Q1内部还自带有一个体二极管（图中箭头所示），它用于防止外部电压反向输入。如果电源方向接反，该二极管会阻断反向电流，防止损坏电路的其他元件。
		
		\item 这个电路的整体工作流程是：当插入电源时，输入电压经过C0滤波以确保其平稳；然后通过S1控制是否将电压VDD传递到下游电路。当S1闭合时，Q1导通，电源得以供给后续电路；当S1断开时，Q1截止，电源被切断，防止任何电流流动。这样设计的电路在实际应用中非常广泛，例如用于便携式电子设备中，用户可以方便地通过机械开关来控制电路是否供电，同时也具备反向电流保护的功能，防止因误操作造成电路的损坏。
	\end{itemize}
	
	\item Analyze the voltage and power amplification of the LM4863 amplifier circuit.
	
	\begin{figure}[htbp]
		\centering
		\includegraphics[width=0.7\textwidth]{ET2_1_analysis.jpg}
		\caption{The analysis of circuit}
		\label{fig:analysis}
	\end{figure}
	
	As we can see in \cref{fig:analysis}, it's exactly the same between the left tone and right tone. Thus, we just need to analyze the half.
	
	For Amp 1L, $v_0=0$, $v_1=-\frac{R_1}{Z_i}v_{in}$, and $\frac{v_1-v_3}{20k\Omega}=\frac{v_3-v_2}{20k\Omega}$, so that $v_3 = (v_1 + v_2) / 2$.
	
	As op amplifier, $v_3=0\Rightarrow v_2=-v_1=\frac{R_1}{Z_i}v_{in}\Rightarrow v_{out} = v_2 - v_1 = \frac{2R_1}{Z_i}v_{in}$.
	
	Thus, voltage gain $A_v$:
	
	\[
	A_v = \frac{v_{\text{out}}}{v_{\text{in}}} = \frac{2R_1}{Z_i} = \frac{2R_1}{R_3 + \frac{1}{j\omega C_3}} = \frac{2j\omega C_3 R_1}{1 + j\omega C_3 R_3}
	\]
	
	\[
	|A_v| = \frac{4 \pi f C_3 R_1}{\sqrt{1 + (2 \pi f C_3 R_3)^2}}, \quad \arg(A_v) = \frac{\pi}{2} - \arctan(2 \pi f C_3 R_3)
	\]
	
	Generally, we can regard $C_3$ only as a filter capacitor ($C_3 \to \infty$) and calculate:
	
	\[
	|A_v| = \frac{2R_1}{R_3} = 8.51, \quad |A_v|_{\text{dB}} = 20 \log \frac{R_1}{R_3} = 18.6 \, \text{dB}
	\]
	
	System cutoff frequency:
	
	\[
	\text{Assume } |A_v| = \frac{\sqrt{2}}{2} \cdot \frac{2R_1}{R_3} \Rightarrow f_L = \frac{1}{2 \pi C_3 R_3} = 33.86 \, \text{Hz}
	\]
	
	We can see that the cutoff frequency is determined by the input RC filter network. That's because $v_0$ is equal to zero in small signal circuit and input signal just generate and propagate circuit instead of voltage into the amplifier.
	
	\[
	P_{\text{out}} = \frac{V_{\text{out}}^2}{R_L}, \quad S_{\text{in}} = \dot{U} \dot{I}^* = v_{\text{in}} \left( v_{\text{in}}^* \cdot \frac{1}{R_3 + j \omega C_3} \right) = \frac{|v_{\text{in}}|^2}{R_3 + j \omega C_3},
	\]
	
	\[
	P_{\text{in}} = \text{Re}(S_{\text{in}}) = \frac{(2\pi f C_3)^2 R_3}{1 + (2\pi f C_3 R_3)^2} |v_{\text{in}}|^2
	\]
	
	\[
	\Rightarrow G_p = \frac{P_{\text{out}}}{P_{\text{in}}} = \frac{1 + (2\pi f C_3 R_3)^2}{(2\pi f C_3)^2 R_3 R_L} \frac{|V_{\text{out}}|^2}{|V_{\text{in}}|^2} \approx 17.0
	\]
	
	\[
	(G_p)_{\text{dB}} = 10 \log G_p = 12.3 \, \text{dB @ } 1 \, \text{kHz}
	\]
	
	The above analysis results can be plotted as shown in the \cref{fig:outcome1} and \cref{fig:outcome2}.
	
	\begin{figure}[htbp]
		\centering
		\includegraphics[width=0.7\textwidth]{ET2_1_analysis_outcome1.jpg}
		\caption{The analysis outcome of circuit ($A_v$-f)}
		\label{fig:outcome1}
	\end{figure}
	
	\begin{figure}[htbp]
		\centering
		\includegraphics[width=0.7\textwidth]{ET2_1_analysis_outcome2.jpg}
		\caption{The analysis outcome of circuit (phase)}
		\label{fig:outcome2}
	\end{figure}
\end{enumerate}

% 讨论
\subsubsection{Discussion}
In this part, we will analyze and discuss the data we got in \cref{tab:tab1} and \cref{tab:tab2}, according to previous analysis.
\begin{enumerate}
	\item 
\end{enumerate}

% 总结
\subsubsection{Conclusion}
\begin{itemize}
	\item 
\end{itemize}
